\documentclass[aps,pra,twocolumn,superscriptaddress,floatfix,longbibliography]{revtex4-2}

\usepackage{amsmath,amssymb}
\usepackage{graphicx}
\usepackage[colorlinks=true,allcolors=blue]{hyperref}

\newcommand{\rd}{\,\mathrm{d}}

\begin{document}

\title{The Stark problem in a Debye plasma: a certified coupled-channel
map of field--screening competition}

\author{C.~D.~Reinoso}
\email{cesardanielrrch@gmail.com}
\affiliation{Facultad de Ciencias, Escuela Superior Polit\'ecnica de
Chimborazo (ESPOCH), Riobamba, Ecuador}

\date{\today}

\begin{abstract}
A hydrogen-like atom in a static electric field and a Debye plasma
experiences two competing ionization mechanisms: field ionization
through the Stark barrier and bound-state dissolution by screening. The
combined problem is non-separable, and the published record consists of
isolated parameter points. We solve it with a coupled-channel
complex-scaled method whose radial channels carry the exact Laguerre
basis of the unperturbed problem, coupled only by the dipole selection
rule $\ell\to\ell\pm1$: the resulting block-tridiagonal complex-symmetric
pencils reach the classical pure-Stark benchmark at $F=0.1$ a.u.\
($E=-0.527418175-0.007269057i$) with fewer than $10^{3}$ degrees of
freedom and geometric convergence in the channel cutoff, and every
assembly is certified at run time against the exact hydrogen spectrum
--- a discipline that caught a factor-of-two kinetic-energy defect
masked by an apparently convergent spurious state during development,
and that directly addresses the documented fragility of
complex-rotated Laguerre-mesh width calculations. On this basis we
compute the first systematic resonance map of the combined problem:
ground-state positions and widths on a grid in the $(F,\lambda)$ plane
and the $m=0$ Stark manifold of $n=2$ under screening. The map
quantifies the field--screening competition and its two regimes, and
the $n=2$ manifold exhibits the interplay of linear Stark splitting
with screening-induced level upshift.
\end{abstract}

\maketitle

\section{Introduction}

The Stark effect in hydrogen and the screened Coulomb (Yukawa--Debye)
potential are two classical single-parameter problems whose resonance
structure is known to high precision: the pure Stark resonances from
complex-dilation and resummation
methods~\cite{Herbst1979,Kolosov1987,Telnov1989,Jentschura2001}, and
the screened-Coulomb bound states and critical screenings from decades
of plasma-physics work~\cite{RogersGraboskeHarwood1970,
delValleNader2018}. Their combination --- an atom in a plasma
subjected to a static field, the generic situation in dense-plasma
spectroscopy --- is qualitatively richer: the field tilts the
potential and opens a Stark barrier, screening lifts and delocalizes
the bound states, and the two ionization mechanisms compete. The
combined problem is non-separable (screening breaks the parabolic
separability of the Stark problem), and the published record consists
of a small number of parameter points: the pioneering ground-state
calculations of Yu and Ho~\cite{YuHo2005}, the complex-scaled
Lagrange-mesh study of Lin and Ho~\cite{LinHo2011} covering ground and
selected excited states, and recent bound-state-level work in dense
plasmas~\cite{DasGhoshal2025}. A systematic map of the
$(F,\lambda)$ plane --- positions, widths, critical curves, and the
fate of the excited Stark manifolds under screening --- has not been
available.

This paper supplies such a map, built on two methodological choices.
First, instead of a two-dimensional mesh, we exploit the dipole
selection rule: the field couples only adjacent angular channels, so
the problem is a block-tridiagonal radial system in which each channel
carries the \emph{exact} orthonormal Laguerre basis of the
unperturbed radial problem~\cite{companionNR}, and the whole pencil is
complex-scaled analytically. Convergence in the channel cutoff is
geometric, and the classical pure-Stark benchmark is reproduced to ten
digits with fewer than a thousand degrees of freedom --- two orders of
magnitude below mesh requirements. Second, every assembly is certified
at run time against exact identities (the unperturbed hydrogen
spectrum, channel overlaps), a discipline inherited from the companion
programs~\cite{companionNR,companionA}. This is not decorative: an
early build of the present code carried a factor-of-two kinetic-energy
defect that produced a \emph{convergent, rotation-stable} spurious
state at a plausible energy --- caught in seconds by the
exact-spectrum certification, not by convergence studies. The
fragility of complex-rotated Laguerre-mesh width calculations against
matrix-element errors is documented in the
literature, and certification is our systematic answer to it.

\section{Method}
\label{sec:method}

In the $m=0$ sector the Hamiltonian of a hydrogen-like atom in a
Debye plasma and a static field $F\hat z$,
\begin{equation}
H=-\tfrac12\nabla^{2}-\frac{Z e^{-\lambda r}}{r}+F z,
\end{equation}
acts on channel wave functions $u_{\ell}(r)$ coupled only by the
dipole matrix elements of $z=r\cos\theta$:
$\langle\ell+1,0|\cos\theta|\ell,0\rangle=
(\ell+1)/\sqrt{(2\ell+1)(2\ell+3)}$. Each channel is expanded in the
orthonormal generalized Laguerre basis of
Ref.~\cite{companionNR}, with a common scale and quadrature grid
across channels so that the inter-channel dipole blocks are exact
contractions. Complex scaling $r\to re^{i\theta_{c}}$ acts
analytically on every block --- kinetic terms scale by
$e^{-2i\theta_{c}}$, Coulombic parts by $e^{-i\theta_{c}}$, the
linear Stark coupling by $e^{+i\theta_{c}}$, and the screening factor
enters the quadrature at complex argument~\cite{AguilarCombes1971,
BalslevCombes1971,Herbst1979,Ho1983,Moiseyev1998} --- giving a
block-tridiagonal complex-symmetric pencil. Resonances are identified
by $\theta_{c}$-stability between two rotation angles, and every
tabulated value carries the maximum spread over a grid of two channel
cutoffs and two basis sizes.

Certification: at $F=\lambda=0$ the assembled pencil must reproduce
the exact hydrogen levels $-1/2n^{2}$; we require agreement below
$10^{-9}$ on the lowest five levels at every build, together with
channel-overlap identities at the $10^{-12}$ level. At $F=0.1$,
$\lambda=0$ the method gives
$E=-0.527418175-0.007269057i$ with geometric convergence in the
channel cutoff ($4\times10^{-5}\to5\times10^{-10}$ over
$\ell_{\max}=4\to16$) at $780$ degrees of freedom, in agreement with
the classical complex-dilation value~\cite{Kolosov1987,Telnov1989,
Jentschura2001}; in the opposite limit $F=0$, $\lambda=0.25$ the
$\ell=1$ block reproduces the independently triple-verified $2p$
resonance of the one-dimensional engine~\cite{companionNR}. A remark
on solver scale: at these dimensions dense complex diagonalization
computes each map point in seconds, so contour-integral eigensolvers,
useful at mesh scales, are unnecessary here.

\section{Results}
\label{sec:results}

[[RESULTS-PLACEHOLDER: D1a anchor table, the (F,lambda) ground map
figure and discussion of the two regimes, the n=2 manifold figure.]]

\section{Conclusions}

[[CONCLUSIONS-PLACEHOLDER]]

\begin{acknowledgments}
The author thanks the ESPOCH Faculty of Sciences for institutional
support. All scripts reproducing the tables and figures, including the
certification batteries, are openly available with the companion
papers.
\end{acknowledgments}

\bibliography{refs}

\end{document}
